\chapter{2014年浙江卷（理）}

\section{单选题}

\begin{problem}
设全集 \(U = \{x \in \mathbf{N} \mid x \geqslant 2\}\)，集合 \(A = \{x \in \mathbf{N} \mid x^2 \geqslant 5\}\)，则 \(\complement_U A =\)（\quad）

\choices
    {\(\varnothing\)}
    {\(\{2\}\)}
    {\(\{5\}\)}
    {\(\{2,5\}\)}
\answer{B}
\end{problem}

\begin{solution}
在全集 \(U\) 中已有 \(x\geqslant2\). 当 \(x=2\) 时，\(x^2=4<5\)，所以 \(2\notin A\)；当 \(x\geqslant3\) 时，\(x^2\geqslant9>5\)，所以 \(x\in A\). 因而 \(U\) 中只有元素 \(2\) 不属于 \(A\)，故 \(\complement_U A=\{2\}\)，选 B.
\end{solution}

\begin{problem}
已知 \(\i\) 是虚数单位，\(a\)，\(b \in \R\)，则 “\(a = b = 1\)” 是 “\((a + b\i)^2 = 2\i\)” 的（\quad）

\choices
    {充分不必要条件}
    {必要不充分条件}
    {充分必要条件}
    {既不充分也不必要条件}
\answer{A}
\end{problem}

\begin{solution}
由复数相等，\((a+b\i)^2=a^2-b^2+2ab\i=2\i\)，所以 \(a^2-b^2=0\)，且 \(ab=1\). 前式给出 \(a=b\) 或 \(a=-b\). 若 \(a=-b\)，则 \(ab=-a^2\leqslant0\)，不可能等于 \(1\)；因此 \(a=b\)，再由 \(ab=1\) 得 \(a=b=1\) 或 \(a=b=-1\).

所以 \(a=b=1\) 能推出 \((a+b\i)^2=2\i\)，但后者还可以由 \(a=b=-1\) 得到，故前者是后者的充分不必要条件，选 A.
\end{solution}

\begin{problem}
某几何体的三视图（单位：\(\mathrm{cm}\)）如图所示. 则此几何体的表面积是（\quad）

\bitmapfigure[max width=0.78\linewidth]{img/2014/zhejiang_science/q3_fig1.png}

\choices
    {\(90 \mathrm{~cm}^{2}\)}
    {\(129 \mathrm{~cm}^{2}\)}
    {\(132 \mathrm{~cm}^{2}\)}
    {\(138 \mathrm{~cm}^{2}\)}
\answer{D}
\end{problem}

\begin{solution}
由三视图可知，该几何体可看作一个底面为直角三角形（两直角边为 \(3\)，\(4\)，斜边为 \(5\)）、棱长为 \(3\) 的直三棱柱，与一个长、宽、高分别为 \(4\)，\(6\)，\(3\) 的长方体拼接而成，拼接面为 \(3\times3\) 的矩形.

直三棱柱的表面积为 \(2\times\frac12\times3\times4+(3+4+5)\times3=48\)，长方体的表面积为 \(2(4\times6+6\times3+4\times3)=108\). 拼接面在两部分表面积中各计算了一次，需减去两次拼接面的面积 \(2\times3\times3=18\). 因而总表面积为 \(48+108-18=138\ \mathrm{cm}^{2}\)，选 D.
\end{solution}

\begin{problem}
为了得到函数 \( y = \sin 3x + \cos 3x \) 的图象，可以将函数 \( y = \sqrt{2}\cos 3x \) 的图象（\quad）

\choices
    {向右平移 \(\frac{\pi}{4}\) 个单位}
    {向左平移 \(\frac{\pi}{4}\) 个单位}
    {向右平移 \(\frac{\pi}{12}\) 个单位}
    {向左平移 \(\frac{\pi}{12}\) 个单位}
\answer{C}
\end{problem}

\begin{solution}
利用和差化积公式，\(\sin3x+\cos3x=\sqrt{2}\cos\left(3x-\frac\pi4\right)=\sqrt{2}\cos3\left(x-\frac\pi{12}\right)\).

因此，函数 \(y=\sin3x+\cos3x\) 的图象由 \(y=\sqrt{2}\cos3x\) 的图象向右平移 \(\frac\pi{12}\) 个单位得到，选 C.
\end{solution}

\begin{problem}
在 \((1 + x)^6 (1 + y)^4\) 的展开式中，记 \(x^{m}y^n\) 项的系数为 \(f(m,n)\)，则 \(f(3,0) + f(2,1) + f(1,2) + f(0,3) =\)（\quad）

\choices
    {45}
    {60}
    {120}
    {210}
\answer{C}
\end{problem}

\begin{solution}
由二项式定理，\(f(3,0)=\mathrm C_6^3=20\)，\(f(2,1)=\mathrm C_6^2\mathrm C_4^1=15\times4=60\)，\(f(1,2)=\mathrm C_6^1\mathrm C_4^2=6\times6=36\)，\(f(0,3)=\mathrm C_6^0\mathrm C_4^3=4\).

所以 \(f(3,0)+f(2,1)+f(1,2)+f(0,3)=20+60+36+4=120\)，选 C.
\end{solution}

\begin{problem}
已知函数 \( f(x) = x^3 + ax^2 + bx + c \)，且 \( 0 < f(-1) = f(-2) = f(-3) \leqslant 3 \)，则（\quad）

\choices
    {\(c \leqslant 3\)}
    {\(3 < c \leqslant 6\)}
    {\(6 < c \leqslant 9\)}
    {\(c > 9\)}
\answer{C}
\end{problem}

\begin{solution}
由 \(f(-1)=f(-2)\)，得 \(-1+a-b+c=-8+4a-2b+c\)，即 \(3a-b=7\). 由 \(f(-2)=f(-3)\)，得 \(-8+4a-2b+c=-27+9a-3b+c\)，即 \(5a-b=19\). 解得 \(a=6\)，\(b=11\).

于是 \(f(-1)=-1+6-11+c=c-6\). 题设给出 \(0<c-6\leqslant3\)，所以 \(6<c\leqslant9\)，选 C.
\end{solution}

\begin{problem}
在同一直角坐标系中，函数 \( f(x) = x^a (x \geqslant 0) \)，\( g(x) = \log_a x \) 的图象可能是（\quad）

\choices
    {\choicebitmap[width=0.15\paperwidth]{img/2014/zhejiang_science/q7_optA.png}}
    {\choicebitmap[width=0.15\paperwidth]{img/2014/zhejiang_science/q7_optB.png}}
    {\choicebitmap[width=0.15\paperwidth]{img/2014/zhejiang_science/q7_optC.png}}
    {\choicebitmap[width=0.15\paperwidth]{img/2014/zhejiang_science/q7_optD.png}}
\answer{D}
\end{problem}

\begin{solution}
由于 \(g(x)=\log_a x\) 有定义，必须有 \(a>0\) 且 \(a\neq1\). 若 \(0<a<1\)，则 \(f(x)=x^a\) 在 \(x\geqslant0\) 上递增且图象向下凹，经过原点和点 \((1,1)\)；同时 \(g(x)=\log_a x\) 递减，经过点 \((1,0)\). 这与选项 D 的两条曲线的单调性、凹凸性和关键点相符.

若 \(a>1\)，则 \(g\) 应递增；因此其余图象均不符合两函数的共同参数关系，选 D.
\end{solution}

\begin{problem}
记 \( \max\{x, y\} = \begin{cases}
x, & x \geqslant y,\\
y, & x < y,
\end{cases} \) \( \min\{x, y\} = \begin{cases}
y, & x \geqslant y,\\
x, & x < y,
\end{cases} \) 设 \(\bs a\)，\(\bs b\) 为平面向量，则（\quad）

\choices
    {\(\min \left\{\left|\bs a + \bs b\right|, \left|\bs a - \bs b\right|\right\} \leqslant \min \left\{\left|\bs a\right|, \left|\bs b\right|\right\}\)}
    {\(\min \left\{\left|\bs a + \bs b\right|, \left|\bs a - \bs b\right|\right\} \geqslant \min \left\{\left|\bs a\right|, \left|\bs b\right|\right\}\)}
    {\(\max \left\{\left|\bs a + \bs b\right|^2, \left|\bs a - \bs b\right|^2\right\} \leqslant \left|\bs a\right|^2 + \left|\bs b\right|^2\)}
    {\(\max \left\{\left|\bs a + \bs b\right|^2, \left|\bs a - \bs b\right|^2\right\} \geqslant \left|\bs a\right|^2 + \left|\bs b\right|^2\)}
\answer{D}
\end{problem}

\begin{solution}
由平行四边形恒等式，\(\left|\bs a+\bs b\right|^2+\left|\bs a-\bs b\right|^2=2\left(\left|\bs a\right|^2+\left|\bs b\right|^2\right)\).

两个数的最大值不小于它们的平均值，所以 \(\max\{\left|\bs a+\bs b\right|^2,\left|\bs a-\bs b\right|^2\}\geqslant\left|\bs a\right|^2+\left|\bs b\right|^2\)，选 D.

选项 A 不恒成立，例如取 \(\bs a=(1,0)\)，\(\bs b=(10,0)\)，则
\(\min\{\left|\bs a+\bs b\right|,\left|\bs a-\bs b\right|\}=9>1=\min\{\left|\bs a\right|,\left|\bs b\right|\}\). 选项 B 不恒成立，例如取 \(\bs a=\bs b\neq\bs0\)，此时左边为 \(0\)，右边为 \(\left|\bs a\right|>0\). 选项 C 也不恒成立，仍取 \(\bs a=\bs b\neq\bs0\)，则左边为 \(4\left|\bs a\right|^2\)，右边为 \(2\left|\bs a\right|^2\). 因此只有 D 恒成立.
\end{solution}

\begin{problem}
已知甲盒中仅有 1 个球且为红球，乙盒中有 \(m\) 个红球和 \(n\) 个蓝球（\(m \geqslant 3\)，\(n \geqslant 3\)），从乙盒中随机抽取 \(i\)（\(i = 1\)，\(2\)）个球放入甲盒中.
放入 \(i\) 个球后，甲盒中含有红球的个数记为 \(\xi_i\)（\(i = 1\)，\(2\)）；放入 \(i\) 个球后，从甲盒中取 1 个球是红球的概率记为 \(p_i\)（\(i = 1\)，\(2\)）.
则（\quad）

\choices
    {\(p_1 > p_2\)，\(E(\xi_1) < E(\xi_2)\)}
    {\(p_1 < p_2\)，\(E(\xi_1) > E(\xi_2)\)}
    {\(p_1 > p_2\)，\(E(\xi_1) > E(\xi_2)\)}
    {\(p_1 < p_2\)，\(E(\xi_1) < E(\xi_2)\)}
\answer{A}
\end{problem}

\begin{solution}
放入一个球后，甲盒共有 \(2\) 个球. 若抽到红球，红球数为 \(2\)；若抽到蓝球，红球数为 \(1\)，故 \(p_1=\frac{m}{m+n}\cdot\frac22+\frac{n}{m+n}\cdot\frac12=\frac{2m+n}{2(m+n)}\).

放入两个球后，甲盒共有 \(3\) 个球. 设放入的两个球中红球数为 \(\eta\)，则每次抽取到红球的概率都是 \(\frac{m}{m+n}\)，故 \(E(\eta)=\frac{2m}{m+n}\). 因而 \(E(\xi_1)=1+\frac{m}{m+n}\)，\(E(\xi_2)=1+\frac{2m}{m+n}\)，从而 \(E(\xi_1)<E(\xi_2)\). 条件于甲盒中红球数为 \(\xi_2\)，从中取出一个球为红球的概率是 \(\frac{\xi_2}{3}\)，所以
\(p_2=E\left(\frac{\xi_2}{3}\right)=\frac{1+\frac{2m}{m+n}}{3}=\frac{3m+n}{3(m+n)}\).

计算得 \(p_1-p_2=\frac{n}{6(m+n)}>0\)，所以 \(p_1>p_2\). 因而选 A.
\end{solution}

\begin{problem}
设函数 \( f_{1}(x) = x^{2} \)，\( f_{2}(x) = 2(x - x^{2}) \)，\( f_{3}(x) = \frac{1}{3} |\sin 2\pi x| \)，\( a_{i} = \frac{i}{99} \)，\(i = 0\)，\(1\)，\(2\)，\(\cdots\)，\(99\). 记 \( I_{k} = |f_{k}(a_{1}) - f_{k}(a_{0})| + |f_{k}(a_{2}) - f_{k}(a_{1})| + \cdots + |f_{k}(a_{99}) - f_{k}(a_{98})| \)，\(k = 1\)，\(2\)，\(3\)，则（\quad）

\choices
    {\(I_{1} < I_{2} < I_{3}\)}
    {\(I_{2} < I_{1} < I_{3}\)}
    {\(I_{1} < I_{3} < I_{2}\)}
    {\(I_{3} < I_{2} < I_{1}\)}
\answer{B}
\end{problem}

\begin{solution}
在 \(0\leqslant x\leqslant1\) 上，\(f_1(x)=x^2\) 递增，因此 \(I_1=f_1(1)-f_1(0)=1\).

函数 \(f_2(x)=2x(1-x)\) 在 \([0,\frac12]\) 上递增、在 \([\frac12,1]\) 上递减. 由于网格点中 \(a_{49}=\frac{49}{99}\)，\(a_{50}=\frac{50}{99}\)，且两者函数值相等，故 \(I_2=2f_2\left(\frac{49}{99}\right)=2\times\frac{2\times49\times50}{99^2}=\frac{9800}{9801}<1\).

函数 \(f_3(x)=\frac13|\sin2\pi x|\) 在各个半周期上先增后减. 网格点中，两个峰值所在的网格点为 \(a_{25}\) 和 \(a_{74}\)，并且 \(f_3(a_{25})=f_3(a_{74})=\frac13\cos\frac\pi{198}\)，\(f_3(a_{49})=f_3(a_{50})=\frac13\sin\frac\pi{99}\).
令 \(\theta=\frac\pi{198}\)，则把网格点按单调段相加，得
\[
I_3=\frac13\cos\theta+\left(\frac13\cos\theta-\frac13\sin2\theta\right)+\left(\frac13\cos\theta-\frac13\sin2\theta\right)+\frac13\cos\theta=\frac23\left(2\cos\theta-\sin2\theta\right)
\]
又 \(0<\theta<\frac1{60}\)，由 \(\cos\theta\geqslant1-\frac{\theta^2}{2}\)、\(\sin2\theta<2\theta\)，有
\[
I_3-1>\frac{1-2\theta^2-4\theta}{3}
>\frac{1-2/60^2-4/60}{3}>0
\]
所以 \(I_3>1\). 因此 \(I_2<I_1<I_3\)，选 B.
\end{solution}

\section{填空题}

\begin{problem}
若某程序框图如图所示，当输入 \(50\) 时，则该程序运行后输出的结果是\fillinblank{}.

\bitmapfigure[max width=0.36\linewidth]{img/2014/zhejiang_science/q11_fig1.png}
\answer{6}
\end{problem}

\begin{solution}
输入 \(n=50\) 后，初始 \(S=0\)，\(i=1\). 按流程逐次计算：
\(S=1\)，\(i=2\)；\(S=4\)，\(i=3\)；\(S=11\)，\(i=4\)；\(S=26\)，\(i=5\)；\(S=57\)，\(i=6\).

此时 \(S=57>50\)，循环结束并输出 \(i=6\)，所以填 \(6\).
\end{solution}

\begin{problem}
随机变量 \(\xi\) 的取值为 0，1，2，若 \(P(\xi = 0) = \frac{1}{5}\)，\(E(\xi) = 1\)，则 \(D(\xi) =\)\fillinblank{}.
\answer{\(\frac{2}{5}\)}
\end{problem}

\begin{solution}
设 \(P(\xi=1)=p_1\)，\(P(\xi=2)=p_2\). 由概率和为 \(1\) 以及期望为 \(1\)，有 \(p_1+p_2=\frac45\)，\(p_1+2p_2=1\).

两式相减得 \(p_2=\frac15\)，从而 \(p_1=\frac35\). 因此 \(E(\xi^2)=p_1+4p_2=\frac35+\frac45=\frac75\)，故 \(D(\xi)=E(\xi^2)-[E(\xi)]^2=\frac75-1=\frac25\).
\end{solution}

\begin{problem}
当实数 \(x\)，\(y\) 满足 \(\begin{cases}
x + 2y - 4 \leqslant 0,\\
x - y - 1 \leqslant 0,\\
x \geqslant 1
\end{cases}\) 时，\(1 \leqslant ax + y \leqslant 4\) 恒成立，则实数 \(a\) 的取值范围是\fillinblank{}.
\answer{\(\left[1,\frac{3}{2}\right]\)}
\end{problem}

\begin{solution}
约束条件等价于 \(x\geqslant1\)，\(y\geqslant x-1\)，\(y\leqslant\frac{4-x}{2}\). 两个关于 \(y\) 的界有交集时，需满足 \(x-1\leqslant\frac{4-x}{2}\)，即 \(x\leqslant2\). 可行域是顶点为 \((1,0)\)，\((1,\frac32)\)，\((2,1)\) 的三角形.

线性函数 \(ax+y\) 在三角形上的最值出现在顶点，三个顶点处的值分别为 \(a\)，\(a+\frac32\)，\(2a+1\). 因而需要同时满足
\(1\leqslant a\leqslant4\)，\(1\leqslant a+\frac32\leqslant4\)，\(1\leqslant2a+1\leqslant4\).
这些不等式分别给出 \(a\in[1,4]\)、\(a\in[-\frac12,\frac52]\)、\(a\in[0,\frac32]\)，取交集得 \(a\in[1,\frac32]\).

所以 \(a\in\left[1,\frac32\right]\).
\end{solution}

\begin{problem}
在 8 张奖券中有一，二，三等奖各 1 张，其余 5 张无奖. 将这 8 张奖券分配给 4 个人，每人 2 张，不同的获奖情况有\fillinblank{}种（用数字作答）.
\answer{60}
\end{problem}

\begin{solution}
把一等奖、二等奖、三等奖分别记为三个不同的奖项. 先选出获得两个奖项的那个人，有 \(4\) 种；再选出获得另一个奖项的人，有 \(3\) 种；最后选择这两个奖项，有 \(\mathrm C_3^2=3\) 种，共有 \(36\) 种.

若三个奖项由三个人分别获得，则有 \(4\times3\times2=24\) 种. 每人最多两张奖券，因此不可能由同一个人获得三个奖项. 所以不同的获奖情况共有 \(36+24=60\) 种.
\end{solution}

\begin{problem}
设函数 \( f(x) = \begin{cases}
x^2 + x, & x < 0,\\
-x^2, & x \geqslant 0.
\end{cases} \) 若 \( f(f(a)) \leqslant 2 \)，则实数 \(a\) 的取值范围是\fillinblank{}.
\answer{\((-\infty,\sqrt{2}]\)}
\end{problem}

\begin{solution}
当 \(a\geqslant0\) 时，\(f(a)=-a^2\leqslant0\)，所以 \(f(f(a))=f(-a^2)=a^4-a^2\). 令 \(u=a^2\geqslant0\)，则 \(f(f(a))\leqslant2\) 等价于 \(u^2-u-2\leqslant0\)，即 \((u-2)(u+1)\leqslant0\). 因此 \(0\leqslant a\leqslant\sqrt2\).

当 \(a<0\) 时，\(f(a)=a^2+a\). 若 \(a<-1\)，则 \(f(a)>0\)，从而 \(f(f(a))=-[f(a)]^2\leqslant0\leqslant2\). 若 \(-1\leqslant a<0\)，则 \(-\frac14\leqslant f(a)\leqslant0\)，于是 \(f(f(a))=f(a)^2+f(a)\leqslant2\). 所以所有 \(a<0\) 都满足条件.

综上，\(a\in(-\infty,\sqrt2]\).
\end{solution}

\begin{problem}
设直线 \(x - 3y + m = 0\)（\(m \neq 0\)）与双曲线 \(\frac{x^2}{a^2} - \frac{y^2}{b^2} = 1\)（\(a > 0\)，\(b > 0\)）的两条渐近线分别交于点 \(A\)，\(B\). 若点 \(P(m, 0)\) 满足 \(|PA| = |PB|\)，则该双曲线的离心率是\fillinblank{}.
\answer{\(\frac{\sqrt5}{2}\)}
\end{problem}

\begin{solution}
设 \(q=\frac ba>0\)，则双曲线的两条渐近线为 \(y=qx\) 与 \(y=-qx\). 直线为 \(x-3y+m=0\)，点 \(P=(m,0)\).

渐近线 \(y=qx\) 与直线的交点为 \(A\left(\frac{m}{3q-1},\frac{qm}{3q-1}\right)\)，渐近线 \(y=-qx\) 与直线的交点为 \(B\left(-\frac{m}{1+3q},\frac{qm}{1+3q}\right)\). 由 \(|PA|=|PB|\)，约去非零因子 \(m^2\)，得
\[
\frac{(2-3q)^2+q^2}{(3q-1)^2}=\frac{(2+3q)^2+q^2}{(1+3q)^2}
\]
由于 \(q=\frac13\) 时直线与渐近线 \(y=qx\) 平行，不能产生题设中的交点 \(A\)，故 \(3q-1\neq0\). 令 \(t=3q>0\)，交叉相乘后为
\[
\left(4-4t+\frac{10}{9}t^2\right)(1+t)^2=\left(4+4t+\frac{10}{9}t^2\right)(t-1)^2
\]
化简得 \(8t(9-4t^2)=0\). 因为 \(t>0\)，所以 \(t=\frac32\)，即 \(q=\frac12\).

双曲线的离心率满足 \(e=\sqrt{1+\frac{b^2}{a^2}}=\sqrt{1+q^2}=\sqrt{1+\frac14}=\frac{\sqrt5}{2}\).
\end{solution}

\begin{problem}
如图，某人在垂直于水平地面 \(ABC\) 的墙前面的点 \(A\) 处进行射击训练. 已知点 \(A\) 到墙面的距离为 \(AB\)，某目标点 \(P\) 沿墙面上的射线 \(CM\) 移动，此人为了准确瞄准目标点 \(P\)，需计算由点 \(A\) 观察点 \(P\) 的仰角 \(\theta\) 的大小. 若 \(AB = 15\ \mathrm{m}\)，\(AC = 25\ \mathrm{m}\)，\(\angle BCM = 30^\circ\)，则 \(\tan\theta\) 的最大值是\fillinblank{}.（仰角 \(\theta\) 为直线 \(AP\) 与平面 \(ABC\) 所成角）

\bitmapfigure[max width=0.42\linewidth]{img/2014/zhejiang_science/q17_fig1.png}
\answer{\(\frac{5\sqrt3}{9}\)}
\end{problem}

\begin{solution}
在地面上取 \(BC\) 方向为水平坐标方向. 由 \(AB\perp BC\)，且 \(AB=15\)，\(AC=25\)，得 \(BC=\sqrt{AC^2-AB^2}=20\).

设目标点 \(P\) 从 \(C\) 沿射线 \(CM\) 移动时，在墙面内沿 \(CB\) 方向的水平位移为 \(u\geqslant0\). 因为 \(\angle BCM=30^\circ\)，目标点的高度为 \(u\tan30^\circ=\frac{u}{\sqrt3}\)，其地面投影到 \(A\) 的距离为 \(\sqrt{15^2+(20-u)^2}\). 因而
\[
\tan\theta=\frac{u/\sqrt3}{\sqrt{15^2+(20-u)^2}}
\]
平方并比较可得
\[
\frac{u^2}{3[225+(20-u)^2]}\leqslant\frac{25}{27}
\]
因为不等式等价于 \(25[225+(20-u)^2]-9u^2=(4u-125)^2\geqslant0\). 当 \(u=\frac{125}{4}\) 时取等号，所以 \(\tan^2\theta\leqslant\frac{25}{27}\). 仰角为非负角，故 \(\tan\theta\) 的最大值为 \(\frac{5\sqrt3}{9}\).
\end{solution}

\section{解答题}

\begin{problem}
在 \(\triangle ABC\) 中，内角 \(A\)，\(B\)，\(C\) 所对的边分别为 \(a\)，\(b\)，\(c\)，已知 \(a \neq b\)，\(c = \sqrt{3}\)，\(\cos^2 A - \cos^2 B = \sqrt{3} \sin A \cos A - \sqrt{3} \sin B \cos B\).

\begin{enumerate}
\item 求角 \(C\) 的大小；
\item 若 \(\sin A = \frac{4}{5}\)，求 \(\triangle ABC\) 的面积.
\end{enumerate}
\end{problem}

\begin{solution}
\begin{enumerate}
\item
由倍角公式，原条件化为
\[
-\sin(A+B)\sin(A-B)=\sqrt3\cos(A+B)\sin(A-B)
\]
由于 \(a\neq b\)，由正弦定理可知 \(A\neq B\)，所以 \(\sin(A-B)\neq0\). 因而
\(\sin(A+B)+\sqrt3\cos(A+B)=0\).

又 \(A+B=\pi-C\in(0,\pi)\)，所以 \(A+B=\frac{2\pi}{3}\)，从而 \(C=\frac\pi3\).

\item
由正弦定理，
\[
\frac{c}{\sin C}=\frac{\sqrt3}{\sqrt3/2}=2
\]
故 \(a=2\sin A=\frac85\). 又 \(B=\frac{2\pi}{3}-A\)，且 \(A<\frac{2\pi}{3}\). 因为 \(\sin A=\frac45\)，只能取 \(\cos A=\frac35\)，于是
\[
\sin B=\sin\left(\frac{2\pi}{3}-A\right)=\frac{\sqrt3}{2}\cos A+\frac12\sin A=\frac{3\sqrt3+4}{10}
\]
所以 \(b=2\sin B=\frac{3\sqrt3+4}{5}\). 三角形面积为
\[
S_{\triangle ABC}=\frac12ab\sin C
=\frac12\cdot\frac85\cdot\frac{3\sqrt3+4}{5}\cdot\frac{\sqrt3}{2}
=\frac{18+8\sqrt3}{25}
\]
\end{enumerate}
\end{solution}

\begin{problem}
已知数列 \(\{a_n\}\) 和 \(\{b_n\}\) 满足 \(a_1a_2a_3\cdots a_n = (\sqrt{2})^{b_n}\)（\(n \in \mathbf{N}^*\)）. 若 \(\{a_n\}\) 为等比数列，且 \(a_1 = 2\)，\(b_3 = 6 + b_2\).

\begin{enumerate}
\item 求 \(a_{n}\) 与 \(b_{n}\)；
\item 设 \(c_{n} = \frac{1}{a_{n}} - \frac{1}{b_{n}}\)（\(n \in \mathbf{N}^*\)）. 记数列 \(\{c_{n}\}\) 的前 \(n\) 项和为 \(S_{n}\).

\begin{enumerate}
\item 求 \(S_{n}\)；
\item 求正整数 \(k\)，使得对任意 \(n \in \mathbf{N}^*\)，均有 \(S_k \geqslant S_n\).
\end{enumerate}
\end{enumerate}
\end{problem}

\begin{solution}
\begin{enumerate}
\item
设等比数列 \(\{a_n\}\) 的公比为 \(q\). 由各项乘积等于正数 \((\sqrt2)^{b_n}\)，且 \(a_1=2>0\)，可知 \(q>0\). 因此
\[
a_1a_2\cdots a_n=2^nq^{\frac{n(n-1)}2}
=2^{\,n+\frac{n(n-1)}2\log_2q}
\]
另一方面，\((\sqrt2)^{b_n}=2^{b_n/2}\)，所以
\[
b_n=2n+n(n-1)\log_2q
\]
代入 \(b_3=6+b_2\)，得
\[
6+6\log_2q=6+\bigl(4+2\log_2q\bigr)
\]
从而 \(\log_2q=1\)，即 \(q=2\). 所以
\(a_n=2^n\)，\(b_n=n(n+1)\).

\item
\begin{enumerate}
\item
于是
\[
S_n=\sum_{j=1}^n\left(\frac1{2^j}-\frac1{j(j+1)}\right)
=\left(1-\frac1{2^n}\right)-\left(1-\frac1{n+1}\right)
=\frac1{n+1}-\frac1{2^n}
\]
\item
再计算相邻两项之差：
\[
S_{n+1}-S_n
=\frac1{2^{n+1}}-\frac1{(n+1)(n+2)}
=\frac{(n+1)(n+2)-2^{n+1}}{2^{n+1}(n+1)(n+2)}
\]
当 \(n=1\)，\(2\)，\(3\) 时，该差分别为正；当 \(n=4\) 时为负. 且 \(2^{n+1}>(n+1)(n+2)\) 在 \(n\geqslant4\) 时可由 \(n=4\) 的 \(32>30\) 归纳得到，因此 \(S_n\) 在 \(n\leqslant4\) 时递增，之后递减. 所以使 \(S_k\) 最大的正整数为 \(k=4\).
\end{enumerate}
\end{enumerate}
\end{solution}

\begin{problem}
如图，在四棱锥 \(A - BCDE\) 中，平面 \(ABC \perp\) 平面 \(BCDE\)，\(\angle CDE = \angle BED = 90^\circ\)，\(AB = CD = 2\)，\(DE = BE = 1\)，\(AC = \sqrt{2}\).

\begin{enumerate}
\item 证明：\(DE \perp\) 平面 \(ACD\)；
\item 求二面角 \(B - AD - E\) 的大小.
\end{enumerate}

\bitmapfigure[max width=0.42\linewidth]{img/2014/zhejiang_science/q20_fig1.png}
\end{problem}

\begin{solution}
\begin{enumerate}
\item
在平面 \(BCDE\) 内，\(CD\perp DE\)，\(BE\perp DE\)，所以 \(CD\parallel BE\). 因而四边形 \(BCDE\) 为直角梯形，且
\[
BC=\sqrt{DE^2+(CD-BE)^2}=\sqrt{1^2+(2-1)^2}=\sqrt2
\]
在三角形 \(ABC\) 中，\(AB^2=4=AC^2+BC^2\)，故 \(AC\perp BC\). 两平面 \(ABC\) 与 \(BCDE\) 的交线为 \(BC\)，且平面互相垂直，所以平面 \(ABC\) 内垂直于 \(BC\) 的直线 \(AC\) 垂直于平面 \(BCDE\). 从而 \(AC\perp CD\)，并且因为 \(AC\) 垂直于平面 \(BCDE\)，还有 \(AC\perp DE\). 又已知 \(CD\perp DE\)，且 \(AC\) 与 \(CD\) 是平面 \(ACD\) 内相交的两条直线，因此 \(DE\perp\) 平面 \(ACD\).

\item
建立空间直角坐标系，取
\(D=(0,0,0)\)，\(E=(1,0,0)\)，\(C=(0,2,0)\)，\(B=(1,1,0)\)，\(A=(0,2,\sqrt2)\).
这些坐标满足题设的长度、垂直关系. 平面 \(ABD\) 和平面 \(EAD\) 的法向量可分别取
\(\bs n_1=\overrightarrow{DA}\times\overrightarrow{DB}=(-\sqrt2,\sqrt2,-2)\)，\(\bs n_2=\overrightarrow{DA}\times\overrightarrow{DE}=(0,\sqrt2,-2)\).
于是
\[
\cos\varphi
=\frac{\bs n_1\cdot\bs n_2}{|\bs n_1||\bs n_2|}
=\frac{6}{2\sqrt2\cdot\sqrt6}
=\frac{\sqrt3}{2}
\]
二面角取其锐角，故 \(\varphi=\frac\pi6\).
\end{enumerate}
\end{solution}

\begin{problem}
如图，设椭圆 \(C: \frac{x^2}{a^2} + \frac{y^2}{b^2} = 1\)（\(a > b > 0\)），动直线 \(l\) 与椭圆 \(C\) 只有一个公共点 \(P\)，且点 \(P\) 在第一象限.

\begin{enumerate}
\item 已知直线 \(l\) 的斜率为 \(k\)，用 \(a\)，\(b\)，\(k\) 表示点 \(P\) 的坐标；
\item 若过原点 \(O\) 的直线 \(l_{1}\) 与 \(l\) 垂直，证明：点 \(P\) 到直线 \(l_{1}\) 的距离的最大值为 \(a - b\).
\end{enumerate}

\bitmapfigure[max width=0.42\linewidth]{img/2014/zhejiang_science/q21_fig1.png}
\end{problem}

\begin{solution}
\begin{enumerate}
\item
椭圆在 \(P(x,y)\) 处的切线斜率为
\[
\frac{dy}{dx}=-\frac{b^2x}{a^2y}
\]
因此 \(k=-\frac{b^2x}{a^2y}\)，即 \(x=-\frac{a^2k}{b^2}y\). 代入椭圆方程，得
\[
\frac{a^2k^2y^2}{b^4}+\frac{y^2}{b^2}=1
\]
所以 \(y=\frac{b^2}{\sqrt{b^2+a^2k^2}}\). 由于 \(P\) 在第一象限，\(y>0\)，再由 \(k<0\) 得
\[
x=-\frac{a^2k}{\sqrt{b^2+a^2k^2}}
\]

\item
直线 \(l_1\) 过原点且与斜率为 \(k\) 的直线垂直，故其方程为 \(x+ky=0\). 记 \(t=|k|>0\)，点 \(P\) 到 \(l_1\) 的距离为
\[
d=\frac{|x+ky|}{\sqrt{1+k^2}}
=\frac{(a^2-b^2)t}{\sqrt{b^2+a^2t^2}\sqrt{1+t^2}}
\]
由
\[
(b^2+a^2t^2)(1+t^2)-(a+b)^2t^2=(at^2-b)^2\geqslant0
\]
得到
\[
d^2\leqslant\frac{(a^2-b^2)^2}{(a+b)^2}=(a-b)^2
\]
当 \(at^2=b\)，即 \(t=\sqrt{\frac ba}\) 时取等号，而且相应切点仍在第一象限. 因此距离的最大值为 \(a-b\).
\end{enumerate}
\end{solution}

\begin{problem}
已知函数 \( f(x) = x^3 + 3|x - a| \)（\(a \in \R\)）.

\begin{enumerate}
\item 若 \( f(x) \) 在 \([-1, 1]\) 上的最大值和最小值分别记为 \(M(a)\)，\(m(a)\)，求 \( M(a) - m(a) \)；
\item 设 \( b \in \R \)，若 \(\left[f(x) + b\right]^2 \leqslant 4\) 对 \( x \in [-1, 1]\) 恒成立，求 \( 3a + b \) 的取值范围.
\end{enumerate}
\end{problem}

\begin{solution}
\begin{enumerate}
\item
当 \(a\leqslant-1\) 时，\(x-a\geqslant0\)，所以
\(f(x)=x^3+3x-3a\) 在 \([-1,1]\) 上递增，故
\(M(a)=4-3a\)，\(m(a)=-4-3a\)，从而 \(M(a)-m(a)=8\).

当 \(a\geqslant1\) 时，\(x-a\leqslant0\)，所以
\(f(x)=x^3-3x+3a\) 在 \([-1,1]\) 上递减，故
\(M(a)=3a+2\)，\(m(a)=3a-2\)，从而 \(M(a)-m(a)=4\).

当 \(-1<a<1\) 时，
\[
f(x)=
\begin{cases}
x^3-3x+3a,&x<a,\\
x^3+3x-3a,&x\geqslant a.
\end{cases}
\]
左段导数为 \(3(x^2-1)\leqslant0\)，右段导数为 \(3(x^2+1)>0\)，所以最小值在 \(x=a\) 处取得，即 \(m(a)=a^3\). 两端点的函数值为 \(f(-1)=3a+2\)，\(f(1)=4-3a\). 它们相等的分界点为 \(a=\frac13\). 因此
\[
M(a)-m(a)=
\begin{cases}
8,&a\leqslant-1,\\
-a^3-3a+4,&-1<a\leqslant\frac13,\\
-a^3+3a+2,&\frac13<a<1,\\
4,&a\geqslant1.
\end{cases}
\]

\item
第二问中，\([f(x)+b]^2\leqslant4\) 等价于 \(-2\leqslant f(x)+b\leqslant2\) 对所有 \(x\in[-1,1]\) 成立. 所以必须且只需
\[
b\in[-2-m(a),\,2-M(a)]
\]
并且 \(M(a)-m(a)\leqslant4\).

由上面的分段式，\(a\leqslant-1\) 时不可能；当 \(-1<a\leqslant\frac13\) 时，条件化为 \(a\geqslant0\)；当 \(\frac13<a<1\) 时恒成立；当 \(a\geqslant1\) 时等号成立. 故 \(a\geqslant0\).

当 \(0\leqslant a\leqslant\frac13\) 时，
\[
3a+b\in[3a-2-a^3,\,6a-2]
\]
当 \(\frac13<a<1\) 时，
\[
3a+b\in[3a-2-a^3,\,0]
\]
当 \(a\geqslant1\) 时 \(b=-3a\)，所以 \(3a+b=0\). 对于 \(0\leqslant a\leqslant\frac13\)，记 \(L(a)=3a-2-a^3\)，\(U(a)=6a-2\). 此时 \(U(a)-L(a)=3a+a^3\geqslant0\)，且 \(L(a)=-2+a(3-a^2)\geqslant-2\)，\(U(a)\leqslant0\)，所以每个区间 \([L(a),U(a)]\) 都包含于 \([-2,0]\). 又 \(U(a)\) 在 \([0,\frac13]\) 上连续递增，且 \(U(0)=-2\)，\(U(\frac13)=0\)，而 \(U(a)\in[L(a),U(a)]\)，故这些区间的并集包含 \([-2,0]\). 当 \(\frac13<a<1\) 时，\(L'(a)=3(1-a^2)>0\)，从而 \(L(a)>L(\frac13)=-\frac{28}{27}>-2\)，且 \(L(a)<L(1)=0\)，所以区间 \([L(a),0]\) 仍包含于 \([-2,0]\). 因而 \(3a+b\) 的取值范围是 \([-2,0]\).
\end{enumerate}
\end{solution}
